\section{Introdução}

Este trabalho trata da detecção e da estimação de sinais impulsivos imersos em
ruído, situação comum em instrumentação e aquisição de dados. A observação em
cada janela de medição é modelada como
\begin{equation}
y[n] = A\, s[n - n_0] + r[n],
\end{equation}
em que \(s[n]\) é a forma conhecida do pulso (template), \(A\) é a amplitude,
\(n_0\) é o instante de início do pulso na janela e \(r[n]\) é o ruído.

O objetivo não é apenas decidir se há pulso, mas também localizar \(n_0\) e
estimar \(A\). Para isso, aplicam-se filtros estocásticos --- em especial o
filtro de Wiener, o filtro casado e o estimador BLUE --- e comparam-se
desempenhos em três tarefas: reconstrução da forma de onda, detecção e
estimação de amplitude.

Os filtros e o limiar de detecção são ajustados apenas no conjunto de treino.
O conjunto de teste é usado uma única vez, após a escolha dos métodos finais.
Os resultados numéricos e as figuras deste relatório foram obtidos a partir da
análise descrita no caderno Jupyter entregue junto com o PDF.
